\documentclass[10.5pt,a4paper]{article}
\usepackage[utf8]{inputenc}
\usepackage[T1]{fontenc}
\renewcommand{\contentsname}{Table des matières}
\usepackage{amsmath,amssymb}
\usepackage{geometry}
\geometry{margin=2cm}
\usepackage{xcolor}
\usepackage{enumitem}
\usepackage{booktabs}
\usepackage{tabularx}
\usepackage{array}
\usepackage{mdframed}
\usepackage{hyperref}
\hypersetup{colorlinks=true,linkcolor=blue!50!black}

\definecolor{ans}{RGB}{20,90,50}
\definecolor{cbleu}{RGB}{30,90,200}

\newmdenv[
  linecolor=cbleu, backgroundcolor=cbleu!5,
  linewidth=1pt, roundcorner=4pt, innermargin=8pt, innertopmargin=6pt, innerbottommargin=6pt,
  skipabove=8pt, skipbelow=8pt
]{objectif}

\newmdenv[
  linecolor=ans, backgroundcolor=ans!6,
  linewidth=1pt, roundcorner=4pt, innermargin=8pt, innertopmargin=6pt, innerbottommargin=6pt,
  skipabove=8pt, skipbelow=8pt
]{test}

\newmdenv[
  linecolor=orange!70!black, backgroundcolor=orange!6,
  linewidth=1pt, roundcorner=4pt, innermargin=8pt, innertopmargin=6pt, innerbottommargin=6pt,
  skipabove=8pt, skipbelow=8pt
]{piege}

% tableau saisie
\newenvironment{saisie}{%
\begin{center}\small
\begin{tabular}{@{}>{\ttfamily}l l@{}}
\toprule
\textrm{\textbf{Le script demande}} & \textbf{Ce que vous tapez} \\
\midrule
}{%
\bottomrule
\end{tabular}
\end{center}
}

\title{\textbf{Scripts Python calculatrice -- Thermodynamique}\\[3pt]
\large Mode d'emploi et exercices de test (Casio)}
\author{}
\date{}

\begin{document}
\maketitle
\vspace{-0.9cm}
\begin{center}
\small\itshape Sept scripts et 79 fiches écran pour traiter les exercices types de l'examen. Tous ont été vérifiés sur les
exercices corrigés du cours (résultats identiques).
\end{center}
\vspace{0.3cm}

\section{Installation et principe général}

\subsection*{Mettre les scripts sur la calculatrice}
\begin{enumerate}[leftmargin=1.5em]
\item Brancher la Casio en USB, ouvrir le logiciel de transfert.
\item Copier les fichiers \texttt{.py} dans la mémoire de la calculatrice.
\item Sur la calculatrice : menu \textbf{Python}, sélectionner le script, \textbf{EXE} pour lancer.
\end{enumerate}

\begin{piege}
\textbf{Trois règles à connaître avant de commencer :}
\begin{itemize}[leftmargin=1.3em,nosep]
\item \textbf{Appuyer sur EXE sans rien taper = « je n'ai pas cette donnée »} $\rightarrow$ le script
saute la question ou la calcule autrement. C'est le mécanisme central : ne tapez que ce que l'énoncé
vous donne.
\item Les \textbf{températures se tapent en $^\circ$C} (le script convertit en Kelvin tout seul), les
\textbf{pressions en bar}, les \textbf{volumes en L}.
\item Le séparateur décimal est le \textbf{point} : tapez \texttt{1.4}, pas \texttt{1,4}.
\end{itemize}
\end{piege}

\subsection*{Les 7 scripts}
\begin{center}
\renewcommand{\arraystretch}{1.35}
\begin{tabularx}{\linewidth}{@{}>{\ttfamily}l X@{}}
\toprule
\textrm{\textbf{Fichier}} & \textbf{À utiliser quand l'exercice parle de...} \\
\midrule
PAC.py & machine frigo, pompe à chaleur, R134a/NH$_3$, diagramme $(\log p,h)$, PSN, COP \\
MOTEUR.py & moteur à piston 4 temps ou 2 temps, essence ou diesel, $\varepsilon$, cylindrée \\
CYCLEGAZ.py & turbine à gaz, cycle de Joule, compresseur + turbine avec un débit \\
TRANSFO.py & une seule transformation d'un gaz parfait (isobare, isochore, isotherme, adiabatique) \\
COMB.py & combustion d'un hydrocarbure, excès d'air $\lambda$, quantité d'air nécessaire \\
THERM.py & échange de chaleur : $Q=mc\Delta T$, chaleur latente, mélange, débit d'un circuit \\
FICHE.py & \textbf{formulaire} en 12 pages plein écran (aucun calcul) : conventions, les 4
transformations, moteurs, combustion, frigo, Joule, vérifications \\
\bottomrule
\end{tabularx}
\end{center}

\subsection*{Les 79 fiches images}
En plus des scripts, le dossier \texttt{images/} contient 79 fiches PNG à la taille exacte de l'écran
($384\times216$ px), à convertir en \texttt{.g3p} avec fx-CG Manager pour les afficher sur la
calculatrice. Le texte est en \textbf{gras et en grande taille} : sur un écran de 384 px, mieux vaut
\emph{plus de fiches} que du texte illisible, donc chaque fiche ne porte qu'une seule idée.
\begin{itemize}[leftmargin=1.4em,nosep]
\item \textbf{D01--D23 (Démarche)} : l'ordre dans lequel rédiger, avec un exemple chiffré en pied de fiche
(essence, diesel, 2 temps, combustion, frigo, PSN/COP, turbine, transformations, échange thermique,
rendements, puissance, vérifications, ordres de grandeur).
\item \textbf{S01--S08 (Schéma)} : les dessins à savoir refaire de tête (cycles $p$-$V$ essence et diesel,
les 4 transformations par un même point, organes du frigo, diagramme $\log p$-$h$, moteur vs récepteur,
fermé vs ouvert, cycle de Joule).
\item \textbf{T01--T35 (Démonstration)} : \textbf{chaque démonstration citée par la prof}, déroulée pas à
pas sur 2 à 4 fiches -- travail élémentaire, 1\textsuperscript{er} principe, machine ouverte, $C_p$ et
$C_v$, adiabatiques, formes des courbes iso-$T$ / adiabatique, machines monothermes, 2\textsuperscript{e}
principe, Carnot, Joule, essence, diesel, entropie, $PSN_{max}$ et $COP_{max}$, PCI/PCS.
\item \textbf{E01--E12 (Exercice résolu)} : les 3 types d'exercices annoncés, entièrement traités
\textbf{avec les nombres} -- cycle de transformations (E01--E03), moteur essence (E04--E08),
frigo/PAC (E09--E12).
\end{itemize}
Elles sont aussi rassemblées et \textbf{expliquées une par une} dans \texttt{fiches-planches.pdf}
(40 pages : chaque fiche avec ce qu'elle montre et le moment où l'ouvrir), et disponibles en JPEG dans
\texttt{images\_jpeg/}.

\newpage
\section{PAC.py -- frigo et pompe à chaleur}

\begin{objectif}
\textbf{Ce qu'il fait :} à partir des enthalpies lues sur le diagramme des frigoristes, il calcule le
travail du compresseur, les chaleurs échangées, le \textbf{PSN} et le \textbf{COP}, puis (en option) les
bornes de Carnot, le débit de fluide frigorigène et le débit du circuit d'eau.
\end{objectif}

\textbf{Ordre des questions :} \texttt{h1} $\rightarrow$ \texttt{h3} $\rightarrow$ \texttt{h2 reel}
(ou \texttt{h2 isos} + \texttt{rend is}) $\rightarrow$ \texttt{Tfroide}/\texttt{Tchaude} $\rightarrow$
\texttt{Putile} + mode $\rightarrow$ circuit secondaire.

\begin{piege}
Si le compresseur n'est \textbf{pas} idéal, laissez \texttt{h2 reel} vide : le script demandera alors
$h_{2s}$ (point isentropique lu sur le diagramme) et $\eta_{isos}$, et calculera $h_2$ réel avec
$h_2=h_1+\dfrac{h_{2s}-h_1}{\eta_{isos}}$. Rappel : $h_4=h_3$ (détente isenthalpique), le script le sait.
\end{piege}

\begin{test}
\textbf{Exercice test 1 -- PAC au R134a.} Évaporation $-5\,^\circ$C, condensation $45\,^\circ$C.
Lecture : $h_1=402{,}5$ ; $h_{2s}=435$ ; $h_3=257$ kJ/kg. $\eta_{isos}=0{,}75$. La maison perd 12 kW.
L'eau des radiateurs se réchauffe de $5\,^\circ$C ($c=4185$ J/kg$\cdot$K).
\emph{Trouver COP, débit de R134a, puissance du compresseur, débit d'eau.}

\begin{saisie}
h1= & 402.5 \\
h3= & 257 \\
h2 reel= & \textrm{(EXE, rien)} \\
h2 isos= & 435 \\
rend is= & 0.75 \\
Tfroide C= & -5 \\
Tchaude C= & 45 \\
Putile kW= & 12 \\
1frigo 2pac: & 2 \\
c sec J/kgK= & 4185 \\
deltaT C= & 5 \\
\end{saisie}

\textbf{Résultats attendus :} \texttt{h2}=445,833 ; \texttt{w}=43,333 ; \texttt{q1}=$-$188,833 ;
\texttt{q2}=145,5 ; \texttt{COP}=\textbf{4,358} ; \texttt{COPmax}=6,363 ;
\texttt{mdot}=\textbf{0,064} kg/s ; \texttt{Pcomp}=\textbf{2,754} kW ; \texttt{mdot2}=\textbf{0,573} kg/s.
\end{test}

\begin{test}
\textbf{Exercice test 2 -- frigo au NH$_3$} (exercices supplémentaires du cours). Lecture sur le
diagramme : $h_1\approx1760$, $h_2\approx2050$, $h_3=h_4\approx700$ kJ/kg. \emph{Trouver la PSN.}

\begin{saisie}
h1= & 1760 \\
h3= & 700 \\
h2 reel= & 2050 \\
Tfroide C= & \textrm{(EXE, rien)} \\
Putile kW= & \textrm{(EXE, rien)} \\
\end{saisie}

\textbf{Résultat attendu :} \texttt{PSN}=\textbf{3,655} \ (soit $\frac{1760-700}{2050-1760}$, la réponse
du corrigé).
\end{test}

\newpage
\section{MOTEUR.py -- moteur à piston}

\begin{objectif}
\textbf{Ce qu'il fait :} il remplit tout le tableau d'état ($T$, $p$, $V$ en A, B, C, D, E), calcule
$W$ et $Q$ de chaque transformation, le rendement théorique, les rendements en cascade, et fait le lien
puissance $\leftrightarrow$ vitesse de rotation. Gère \textbf{essence et diesel}, \textbf{4 temps et 2
temps}.
\end{objectif}

\textbf{Ordre des questions :} type (1 essence / 2 diesel) $\rightarrow$ nombre de temps $\rightarrow$
$\varepsilon$, $\gamma$ $\rightarrow$ conditions d'admission $\rightarrow$ cylindrée totale et nombre de
cylindres $\rightarrow$ données de combustion $\rightarrow$ rendements $\rightarrow$ $N$ ou $P_{eff}$.

\begin{piege}
Pour la combustion, trois façons de renseigner le point D, dans cet ordre de priorité :
\begin{enumerate}[leftmargin=1.4em,nosep]
\item \texttt{QCD J} directement si l'énoncé donne la chaleur ;
\item sinon, en \textbf{diesel} seulement, le \texttt{rho} (rapport de combustion $V_D/V_C$) ;
\item sinon, les données du carburant (\texttt{PCI}, \texttt{x}, \texttt{y} de C$_x$H$_y$,
\texttt{lambda}) et le script calcule lui-même la masse de carburant admise et $Q_{CD}$.
\end{enumerate}
Laissez vide ce que vous n'avez pas. Les rendements laissés vides valent 1.
\end{piege}

\begin{test}
\textbf{Exercice test 3 -- moteur essence} (exercices supplémentaires du cours). 4 temps, 4 cylindres,
cylindrée 1,2 L, $\varepsilon=10$, $P_{eff}=20$ kW. Air admis à $16\,^\circ$C, 1 bar, $\gamma=1{,}4$.
Carburant CH$_{1,8}$, $\lambda=1{,}2$, PCI $=42500$ kJ/kg. $\eta_{forme}=0{,}75$, $\eta_{méca}=0{,}8$.
\emph{Trouver le rendement théorique et la vitesse de rotation.}

\begin{saisie}
1ess 2die: & 1 \\
1=4tps 2=2tps: & 1 \\
eps= & 10 \\
gamma= & 1.4 \\
Tadm C= & 16 \\
padm bar= & 1 \\
cyl tot L= & 1.2 \\
nb cyl= & 4 \\
QCD J= & \textrm{(EXE, rien)} \\
PCI kJ/kg= & 42500 \\
x de CxHy= & 1 \\
y de CxHy= & 1.8 \\
lambda= & 1.2 \\
rend forme= & 0.75 \\
rend meca= & 0.8 \\
N tr/min= & \textrm{(EXE, rien)} \\
Peff kW= & 20 \\
\end{saisie}

\textbf{Résultats attendus :} \texttt{TC}=726,31 K ; \texttt{pC}=25,12 bar ; \texttt{QCD}=875,79 J ;
\texttt{eta th}=\textbf{60,19 \%} (valeur exacte du corrigé) ; \texttt{Wth}=527,13 J ;
\texttt{Wdisp}=316,28 J ; \texttt{N}=\textbf{1897} tr/min.

\emph{Remarque : le corrigé du cours annonce 1881 tr/min. L'écart de 0,8 \% vient uniquement des
arrondis intermédiaires de la correction ; le script, lui, boucle exactement
($Q_{CD}+Q_{EB}+W_{cycle}=0$).}
\end{test}

\begin{test}
\textbf{Exercice test 4 -- moteur diesel.} 4 temps, 4 cylindres, cylindrée 2 L, $\varepsilon=18$,
$\rho=2$, air admis à $27\,^\circ$C et 1 bar, $\gamma=1{,}4$, $N=2200$ tr/min.
\emph{Trouver tous les points, le rendement et la puissance théorique.}

\begin{saisie}
1ess 2die: & 2 \\
1=4tps 2=2tps: & 1 \\
eps= & 18 \\
gamma= & 1.4 \\
Tadm C= & 27 \\
padm bar= & 1 \\
cyl tot L= & 2 \\
nb cyl= & 4 \\
QCD J= & \textrm{(EXE, rien)} \\
rho= & 2 \\
rend forme= & \textrm{(EXE, rien)} \\
rend meca= & \textrm{(EXE, rien)} \\
N tr/min= & 2200 \\
\end{saisie}

\textbf{Résultats attendus :} \texttt{VA}=29,41 cm$^3$ ; \texttt{TC}=953,78 K ; \texttt{pC}=57,2 bar ;
\texttt{TD}=1907,56 K ; \texttt{TE}=792,1 K ; \texttt{pE}=2,64 bar ;
\texttt{eta th}=\textbf{63,16 \%} ; \texttt{QCD}=588,8 J ; \texttt{Wth}=371,88 J ;
\texttt{Pth}=\textbf{27,27 kW}.
\end{test}

\begin{piege}
\textbf{Pour un 2 temps}, répondez \texttt{2} à la question du nombre de temps : le script utilise alors
$P=\frac{N}{60}\,W\,z$ au lieu de $\frac{N}{120}\,W\,z$. C'est la seule différence de calcul (le cycle
thermodynamique est identique à l'essence).
\end{piege}

\newpage
\section{CYCLEGAZ.py -- turbine à gaz, cycle de Joule}

\begin{objectif}
\textbf{Ce qu'il fait :} vous enchaînez les composants un par un (compresseur, chambre de combustion,
turbine, refroidisseur) et le script suit l'état du gaz ($T$, $p$) tout au long, accumule le travail et
la chaleur, puis donne le rendement et les puissances via le débit.
\end{objectif}

\textbf{À chaque étape}, il propose : \texttt{1} adiabatique (compresseur ou turbine) ou \texttt{2}
isobare (échangeur / chambre de combustion), \texttt{3} pour terminer. Pour une adiabatique, deux
possibilités : donner le \textbf{rapport de pression} (\texttt{1}) ou un \textbf{travail imposé}
(\texttt{2}, cas d'une turbine qui entraîne exactement le compresseur).

\begin{test}
\textbf{Exercice test 5 -- turbine à gaz à double détente} (exercice 1 des exercices supplémentaires).
Air, $c_p=1004$ J/kg$\cdot$K, $\gamma=1{,}4$, débit 35 kg/s. 1-2 : compression adiabatique de 1 à 4,1 atm,
$T_1=15\,^\circ$C. 2-3 : isobare jusqu'à $750\,^\circ$C. 3-4 : détente adiabatique dans une turbine dont
tout le travail sert à entraîner le compresseur. 4-5 : isobare jusqu'à $700\,^\circ$C. 5-6 : détente
adiabatique jusqu'à 1 atm. 6-1 : refroidissement isobare.
\emph{Trouver $T_2,T_4,T_6$, $p_4$, le rendement et la puissance utile.}

\begin{saisie}
cp J/kgK= & 1004 \\
gamma= & 1.4 \\
debit kg/s= & 35 \\
T1 C= & 15 \\
p1 (bar/atm)= & 1 \\
\textrm{étape 1 :} 3=fin: & 1 \ \textrm{puis} \ 1 \ \textrm{puis} \ \texttt{p fin=}4.1 \\
\textrm{étape 2 :} 3=fin: & 2 \ \textrm{puis} \ \texttt{T fin C=}750 \\
\textrm{étape 3 :} 3=fin: & 1 \ \textrm{puis} \ 2 \ \textrm{puis} \ \texttt{w J/kg=}-143649 \\
\textrm{étape 4 :} 3=fin: & 2 \ \textrm{puis} \ \texttt{T fin C=}700 \\
\textrm{étape 5 :} 3=fin: & 1 \ \textrm{puis} \ 1 \ \textrm{puis} \ \texttt{p fin=}1 \\
\textrm{étape 6 :} 3=fin: & 2 \ \textrm{puis} \ \texttt{T fin C=}15 \\
\textrm{étape 7 :} 3=fin: & 3 \\
\end{saisie}

\textbf{Résultats attendus :} $T_2=$\textbf{431,22 K} ; $T_4=$\textbf{880,07 K} ; $p_4=$\textbf{2,42} ;
$T_6=$\textbf{756,0 K} ; \texttt{eta}=\textbf{31,7 \%} ; \texttt{Pnet}=\textbf{$-$7631 kW} $=-7{,}63$ MW.
\emph{(Le corrigé donne 431 K, 880 K, 756 K, 31,6 \%, $-7{,}63$ MW.)}
\end{test}

\begin{piege}
À l'étape 3, le travail imposé est celui du compresseur \textbf{changé de signe} : le script a affiché
\texttt{w=143643} à l'étape 1, on entre donc $-143649$ (arrondi sans importance). C'est la traduction de
« tout le travail de la turbine sert à faire fonctionner le compresseur ».
\end{piege}

\newpage
\section{TRANSFO.py -- une transformation de gaz parfait}

\begin{objectif}
\textbf{Ce qu'il fait :} pour \textbf{une seule} transformation en système fermé, il complète l'état
initial manquant, calcule l'état final, puis $W$, $Q$, $\Delta U$ et $\Delta H$. C'est l'outil à
dégainer quand un exercice demande « calculez le travail et la chaleur de la transformation A$\to$B ».
\end{objectif}

\textbf{État initial :} tapez \texttt{p1}, \texttt{V1}, \texttt{T1} et \texttt{n} en laissant
\textbf{une seule} de ces cases vide -- le script la calcule avec $pV=nRT$.
\textbf{État final :} ne remplissez que la grandeur que l'énoncé donne, laissez les deux autres vides.

\begin{test}
\textbf{Exercice test 6 -- compression adiabatique.} Un cylindre contient 0,333 L d'air à 1 bar et
$16\,^\circ$C ($\gamma=1{,}4$). On le comprime de façon adiabatique réversible jusqu'à 0,0333 L
($\varepsilon=10$). \emph{Trouver $T_2$, $p_2$, $W$ et $Q$.}

\begin{saisie}
gamma= & 1.4 \\
p1 bar= & 1 \\
V1 L= & 0.33333 \\
T1 C= & 16 \\
n mol= & \textrm{(EXE, rien)} \\
3isoT 4adia: & 4 \\
p2 bar= & \textrm{(EXE, rien)} \\
V2 L= & 0.033333 \\
T2 C= & \textrm{(EXE, rien)} \\
\end{saisie}

\textbf{Résultats attendus :} \texttt{n}=0,014 mol ; \texttt{T2}=\textbf{726,312 K} ;
\texttt{p2}=\textbf{25,119 bar} ; \texttt{W}=\textbf{+125,989 J} ; \texttt{Q}=\textbf{0} ;
\texttt{dU}=125,989 J.
\emph{(Ce sont exactement les valeurs du point C du moteur essence de l'exercice test 3.)}
\end{test}

\section{COMB.py -- combustion et excès d'air}

\begin{objectif}
\textbf{Ce qu'il fait :} pour un carburant C$_x$H$_y$, il équilibre la combustion, donne les moles
d'O$_2$ et d'air (stœchiométrique et réel avec $\lambda$), le rapport air/carburant en masse, et pour
une quantité donnée de carburant : les moles et la masse d'air nécessaires, plus la chaleur dégagée.
\end{objectif}

\begin{test}
\textbf{Exercice test 7 -- exemple du syllabus.} « J'ai du CH$_{1,9}$ et de l'air (21 \% d'O$_2$,
$M=28{,}9$ g/mol). Quelle quantité d'air me faut-il pour brûler exactement 1 kg de carburant ? »

\begin{saisie}
x= & 1 \\
y= & 1.9 \\
lambda= & 1 \\
m carb kg= & 1 \\
PCI kJ/kg= & \textrm{(EXE, rien)} \\
\end{saisie}

\textbf{Résultats attendus :} \texttt{nO2 stoe}=\textbf{1,475} ; \texttt{nair stoe}=\textbf{7,024} ;
\texttt{M carb}=\textbf{13,9} g/mol ; \texttt{n carb}=\textbf{71,942} mol ;
\texttt{n air}=\textbf{505,31} mol ; \texttt{m air}=\textbf{14,603} kg.
\emph{(Le syllabus trouve 1,475 -- 7,02 -- 13,9 -- 71,94 -- 505,04 -- 14595 g : identique aux arrondis
près.)}
\end{test}

\section{THERM.py -- échanges thermiques}

\begin{objectif}
\textbf{Ce qu'il fait :} un menu à 4 entrées pour toute la partie « échange de chaleur » qui accompagne
souvent l'exercice frigo/chauffage : \texttt{1} chauffage simple $Q=mc\Delta T$, \texttt{2} changement
d'état $Q=mL$, \texttt{3} température finale d'un mélange, \texttt{4} circuit à débit
$P=\dot m\,c\,\Delta T$.
\end{objectif}

\begin{test}
\textbf{Exercice test 8 -- circuit de chauffage.} Une PAC fournit 12 kW à un circuit d'eau qui se
réchauffe de $5\,^\circ$C ($c_{eau}=4185$ J/kg$\cdot$K). \emph{Trouver le débit d'eau.}

\begin{saisie}
4 debit: & 4 \\
P kW= & 12 \\
mdot kg/s= & \textrm{(EXE, rien)} \\
c J/kgK= & 4185 \\
dT C= & 5 \\
\end{saisie}

\textbf{Résultats attendus :} \texttt{mdot}=\textbf{0,573 kg/s} ; \texttt{mdot m3/h}=\textbf{2,065}
m$^3$/h. \emph{(Le script devine ce qu'il faut calculer : laissez vide l'inconnue parmi $P$, $\dot m$ et
$\Delta T$.)}
\end{test}

\vspace{0.5cm}
\begin{piege}
\textbf{Stratégie le jour de l'examen.} Les scripts donnent les \emph{nombres}, mais un exercice se note
sur le \emph{raisonnement} : écrivez toujours la formule utilisée et le type de chaque transformation
avant de donner la valeur. Servez-vous en surtout pour (1) aller vite sur les calculs répétitifs,
(2) \textbf{vérifier} un résultat obtenu à la main, (3) ne pas perdre de points sur une erreur de
calculatrice.
\end{piege}

\end{document}
